\SourceLeaf{328}%
\par\noindent\hspace*{\parindent} Bei der obigen Entwicklung wurde für \(Z\) nur vorausgesetzt, dass es eine von \(u\) abhängige veränderliche Größe ist, deren Ableitung \(\frac{\mathrm{d}Z}{\mathrm{d}u}=z\) durch irgendeine Funktion von \(X\) und \(Y\) ausgedrückt wird. Deshalb kann man die gefundenen Ergebnisse unmittelbar auf jede andere veränderliche Größe übertragen, für die dieselben Voraussetzungen gelten, insbesondere auf \(X\) oder \(Y\) selbst. Man muss nur anstelle von \(L,\mathrm{L},l\) und den verschieden markierten Formen von \(\mathrm{h}\) entsprechend \(T,\mathrm{T},t\) und die verschiedenen Formen von \(\mathrm{f}\) oder \(B,\mathrm{B},b\) und die verschiedenen Formen von \(\mathrm{g}\) einsetzen. Gleichung (4), von der hier sonst kein direkter Gebrauch gemacht wird, liefert zunächst die folgenden beiden Gleichungen, die ebenfalls bis zur vierten Ordnung ausschließlich genau sind.
\[
\begin{aligned}
T&=\mathrm{T}+\tfrac18(\mathrm{ff}'+\mathrm{f}''\mathrm{g})rr\\
B&=\mathrm{B}+\tfrac18(\mathrm{fg}'+\mathrm{gg}'')rr
\end{aligned}
\]
Daraus folgt zunächst, dass sich \(h\) und \(\mathrm{h}\) als Werte von \(z\), je nachdem, ob man \(T\) und \(B\) oder \(\mathrm{T}\) und \(\mathrm{B}\) für \(X\) und \(Y\) einsetzt, um eine Größe zweiter Ordnung unterscheiden. Dieser Unterschied wird bis zur vierten Ordnung ausschließlich genau durch die Formel bestimmt
\[
h=\mathrm{h}+\tfrac18(\mathrm{ff}'+\mathrm{f}''\mathrm{g})rr\mathord{.}(\tfrac{\mathrm{d}z}{\mathrm{d}X})+\tfrac18(\mathrm{fg}'+\mathrm{gg}'')rr(\tfrac{\mathrm{d}z}{\mathrm{d}Y})
\]
wobei für die partiellen Ableitungen \((\frac{\mathrm{d}z}{\mathrm{d}X})\) und \((\frac{\mathrm{d}z}{\mathrm{d}Y})\) oder \(z',z''\) ihre bestimmten Werte bei \(X=\mathrm{T},\ Y=\mathrm{B}\) einzusetzen sind, nämlich \(\mathrm{h}'\) und \(\mathrm{h}''\). Es gilt also mit Genauigkeit bis zur vierten Ordnung
\[
\mathrm{h}=h-\tfrac18(\mathrm{ff'h}'+\mathrm{fg'h}''+\mathrm{f''gh}'+\mathrm{gg''h}'')rr
\]
und durch Einsetzen dieses Werts in Gleichung (5) erhält man mit Genauigkeit bis zur fünften Ordnung ausschließlich
\[
\begin{aligned}
l={}&-hr+\tfrac1{24}(2\mathrm{ff'h}'+2\mathrm{fg'h}''+2\mathrm{f''gh}'+2\mathrm{gg''h}''\\
&\qquad-\mathrm{ffh}'''-2\mathrm{fgh}''''-\mathrm{ggh}^{\mathrm{v}})r^3
\end{aligned}
\]

\indent Aus denselben Gründen, aus denen sich \(\mathrm{h}\) von \(h\) unterscheidet, werden sich auch \(\mathrm{f},\mathrm{f}',\mathrm{f}''\) usw., \(\mathrm{g},\mathrm{g}',\mathrm{g}''\) usw. und \(\mathrm{h}',\mathrm{h}''\) usw. von \(f,f',f''\) usw., \(g,g',g''\) usw. und \(h',h''\) usw. jeweils um Größen zweiter Ordnung unterscheiden. Im gerade angegebenen Ausdruck für \(l\) kann man daher die ersteren Größen durch die letzteren ersetzen, ohne die Genauigkeit zu verringern. Es gilt also ebenfalls mit Genauigkeit bis zur fünften Ordnung ausschließlich
